\RequirePackage{etex}
\newcommand{\DIRClass}{}
\documentclass
[pauta]
%[mayor]
%[menor]
%[basica]
{\DIRClass cmat_grupales_FICHA}
\usepackage{cmat_grupales_MACROS}
\usepackage{graphicx}
\usepackage{tikz}
\usetikzlibrary{shapes.geometric}
\usetikzlibrary{positioning}
\usetikzlibrary{calc}
\usepackage{epstopdf}
\usepackage{subcaption}
\renewcommand{\thesubfigure}{}
\captionsetup[subfigure]{labelformat=empty}

%\newcommand*{\Hex}{draw, regular polygon, regular polygon sides = 6, inner sep=0.25cm}

\newcommand*{\cultivo}{\begin{center}\begin{tikzpicture}[node distance={10.2mm}, main/.style = {draw, regular polygon, regular polygon sides = 6, minimum size=12mm, inner sep=0pt, very thick}]}
\newcommand*{\fin}{\end{tikzpicture}\end{center}}
\newcommand*{\finfin}{\end{tikzpicture}}
\newcommand*{\cultivos}{\begin{tikzpicture}[node distance={10.2mm}, main/.style = {draw, regular polygon, regular polygon sides = 6, minimum size=12mm, inner sep=0pt, very thick}]}
\newcommand*{\granja}{\node[main]}
\newcommand*{\cultivoP}{\begin{center}\begin{tikzpicture}[node distance={7.8mm}, main/.style = {draw, regular polygon, regular polygon sides = 6, minimum size=8.8mm, inner sep=0pt, very thick}]}
\newcommand*{\cultivosP}{\begin{tikzpicture}[node distance={7.8mm}, main/.style = {draw, regular polygon, regular polygon sides = 6, minimum size=8.8mm, inner sep=0pt, very thick}]}
\newcommand*{\granjaN}{\node[main, fill=gray]}
\newcommand*{\grande}{10.2mm}
\newcommand*{\enano}{7.8mm}

%\epstopdfsetup{outdir=./}
%% compilar siempre TeXify -> DVIPS -> ps2pdf
%% en winedt, Shift + Control + X -> Shift + Control + D -> ps2pdf

%%%%%%%%%%%%%%%%%%%%%%%%%%%%%%%%%%%%
%%%%%%%%%%%%%%%%%%%%%%%%%%%%%%%%%%%%
%%%%%                          %%%%%
%%%%%     SECTOR EDITABLE      %%%%%
%%%%%                          %%%%%
%%%%%%%%%%%%%%%%%%%%%%%%%%%%%%%%%%%%
%%%%%%%%%%%%%%%%%%%%%%%%%%%%%%%%%%%%

%% NOTA: en caso de fallas, revisar
%%       'set-up' (linea 66)


%% Propiedades básicas
%% Ajustar qué fecha es (1-5 tradicionalmente), versión, el día de la prueba
%% Y las horas!


\newcommand{\nfecha}{Fecha X, Borrador}% <Ordinal> fecha
\newcommand{\version}{XXIV}% 2003 -> I
\newcommand{\versionbasica}{XXIII}% 2004 -> I
\newcommand{\fecha}{día a definir / 2026}% DD de <mes> de AAAA
\newcommand{\titgrupal}{Colonización Planetaria}
\newcommand{\rutaprueba}{cmat_grupales_TEXTO.tex}


%% NOTA: elegir al final del documento una tabla de puntaje para la pauta



%%
%% Horas: Cambiar según corresponda
%%


\textoalt{
\newcommand{\horainicial}{16:00 }
\newcommand{\horafinal}{18:00}
\newcommand{\horasubida}{18:20}
%Básica
}{
\newcommand{\horainicial}{15:45 }
\newcommand{\horafinal}{17:45}
\newcommand{\horasubida}{18:05}
%Menor
}{
\newcommand{\horainicial}{16:15 }
\newcommand{\horafinal}{18:15}
\newcommand{\horasubida}{18:35}
%Mayor
}{}






%%%%%%%%%%%%%%%%
%%%% Set-up %%%%  NO TOCAR (salvo sea necesario)
%%%%%%%%%%%%%%%%



%%%%%%%%%%%%%%%%%%%%%%%%%%%%%%%%%%%%%%%
%% Para el diseño de pagina
%%

%\geometry{letterpaper, top=1.5cm, bottom=1.2cm, right=2cm, left=2cm, includefoot} % Márgenes de página

\rhead{}
\lhead{}
\rfoot{}% comentar(%) para obtener 'justifique su respuesta'



%%%%%%%%%%%%%%%%%%%%%%%%%%%%%%%
%% En caso de emergencia...
%%

%% para cambios locales
%%
%\psset{xunit=1.0cm,yunit=1.0cm,algebraic=true,dotstyle=o,dotsize=3pt 0,linewidth=0.8pt,arrowsize=3pt 2,arrowinset=0.25}
%\useackage[T1]{fontenc}% codificacion fuentes
%\usepackage[activeacute,spanish]{babel}% idioma
%\usepackage[utf8]{inputenc}% codificacion caracteres
%\usepackage[dvips]{graphicx}% driver colores
\DeclareGraphicsExtensions{.bmp,.png,.pdf,.jpg}

%% si la clase 'falla'
%%
%\usepackage[<option>]{../CMAT-macrosgrupales}

%\puntajefalse
%\basicatrue
%\menortrue
%\mayortrue
%\pautatrue



%%%%%%%%%%%%%%%%%%%%%%%%%%%%%%%%%%%%%%%%%%%%%%%%
%% Otros paquetes posiblemente solicitados...
%%

%\usepackage{wrapfig}% env: wrapfig (mini figuras)
%\usepackage{subfigure}% extencion figure + env: subfigure
%\usepackage[rightcaption]{sidecap}% env: SCfigure (mover caption)

%\usepackage{tabularx}% env: tabularx
%\usepackage{longtable}% env: longtable

%\usepackage{fancyvrb}% cmd: verbatin (cajas de codigo latex)
%\usepackage{hhline}% cmd: hhline
%\usepackage{arcs}% cmd: overarc, underarc
%\usepackage{physics}% +comandos y simbolos

%\usepackage{enumerate}% enumitem cargado por defecto

%\usepackage{hyperref}% NO MOVER, SOLO COMENTAR (%)


%%%%%%%%%%%%%%%%%%%%%%%%%%%%%%%%%%%
%% Posibles comandos comunes...

\providecommand{\abs}[1]{\lvert#1\rvert}
\newcommand{\pares}[1]{\left(#1\right)}
\newcommand{\nariz}[1]{\left\{#1\right\}}

%%%%%%%%%%%%%%%%%%%%%%%%%%%%%%%%%%%%%%%%%%%%%%%%%%%%%%%%%%%%%%%%




\begin{document}


%% NO TOCAR
%% ESTO DEBERÍA FUNCIONAR SOLO
%% (espero)

\author{ }

\begin{minipage}{0.3\textwidth}
\psscalebox{1.6}{
        \psset{unit=0.5cm,algebraic=true,dotstyle=*,dotsize=3pt 0,linewidth=0.8pt,arrowsize=3pt 2,arrowinset=0.25}
        \begin{pspicture}(-5,0)(1,3.2)
        \rput[tl](-5.5,5.2){\includegraphics[width=3cm, bb = 0 0 556 556]{\DIRClass cmat_logo.eps}        }
	    \uput[ur](-1.5,0.315){\psscalebox{3.8}{\textoalt{\versionbasica}{\version}{\version}{\version}}}
        \end{pspicture}
        }\par\bigskip

\end{minipage}\hfill
\begin{minipage}{0.5\textwidth}
%\vspace{-1.2cm}
\begin{flushright}
  {\LARGE\textbf{\nfecha}}\\
  {\fecha}\\
  {\LARGE \textbf{\textoalt{Básica}{Menor}{Mayor}{Borrador Confidencial}}}\\
  {Grupal}
  \end{flushright}
\end{minipage}

%\begin{minipage}{0.3\textwidth}
       % \centering
      %  \includegraphics[scale=0.05]{cmat_logo.pdf}  % Cambia a PDF y ajusta el tamaño según sea necesario
     %   \par\bigskip
    %\end{minipage}\hfill
    %\begin{minipage}{0.5\textwidth}
       % \begin{flushright}
      %      {\LARGE\textbf{\nfecha}}\\
     %       {\fecha}\\
    %        {\LARGE \textbf{\textoalt{Básica}{Menor}{Mayor}{Soluciones}}}\\
   %         {Grupal}
  %      \end{flushright}
 %\end{minipage}

\textoocultoSN{S}{S}{S}{N}{\psset{unit=0.8cm,algebraic=true,dotstyle=*,dotsize=3pt 0,linewidth=0.8pt,arrowsize=3pt 2,arrowinset=0.25}
        \begin{pspicture*}(0,-2.7)(\linewidth,-1.9)
        \psline[linewidth=1pt](0,-1.9)(\linewidth,-1.9)
        \rput[t](0.5\linewidth,-1.9){\fbox{\large\textbf{Entreguen solo una respuesta por equipo.}}}
        \end{pspicture*}}


%%
%% cuadro texto formato online
%%


%\textoocultoSN{S}{S}{S}{N} % visible
\textoocultoSN{N}{N}{N}{N} % oculto
{\fbox{\parbox{\textwidth-6.8pt}{\textbf{Aviso importante:} Tienen \textbf{dos horas} para resolver la prueba, entre las \horainicial y \horafinal. Posterior a ello, tendrán \textbf{sólo} 20 minutos para subir la prueba a la plataforma, hasta las \horasubida. \textbf{Pruebas subidas más tarde, o no subidas a la plataforma, no serán evaluadas}. \\ Recomendamos además escribir no más de tres respuestas por plana, de modo que sea lo más ordenado posible. }}}



%%%%%%%%%%%%%%%%%%%%%%%%%%%%%%%%%%%%%%%
%% ACÁ EMPIEZA EL TEXTO DE LA PRUEBA %%
%%%%%%%%%%%%%%%%%%%%%%%%%%%%%%%%%%%%%%%

\vspace{1em}
\begin{center} {\Large\textbf{\titgrupal}}\end{center}
\vspace{1em}


\newcommand{\saltoextra}{}
\input{\rutaprueba}



%%%%%%%%%%%%%%%%%%%%%%%
%% TABLA DE PUNTAJES %%
%%%%%%%%%%%%%%%%%%%%%%%



\begin{parrafoocultoSN}{N}{N}{N}{S}
\newpage
\section*{Preguntas por nivel}

%% Favor de elegir una sola de las siguientes tablas de acuerdo al interes y dejar el resto comentado (%).

%% descripciones:

% tablapuntaje: Tabla usual. Se usa por defecto.
% Los puntajes de subpreguntas se muestran acompañados por un indice "(<#indice>)" con "+" como separadores.

% tablapuntajeH: Tabla recomendada para pruebas donde hayan muchas preguntas con una cantidad mas menos consistente de subpreguntas.
% Se muestran todos los puntajes de las subpreguntas usando "+" como separadores.

% tablapuntajeV: Tabla recomendada para pruebas donde hayan pocas preguntas con muchas subpreguntas.
% Los puntajes de subpreguntas se muestran en pisos separados.

\tablapuntaje
%\tablapuntajeH
%\tablapuntajeV

\end{parrafoocultoSN}




\end{document} 
